%=========================================================================
%  INHALT — layoutneutral.
%
%  Hier steht ausschliesslich, WAS im Lebenslauf steht. WIE es aussieht,
%  entscheidet das Theme: die hier verwendeten Makros (\cvsection, \cvrole,
%  cvbullets, \cvfact, \cvtag, \lead) definiert jedes Theme selbst.
%  Ein Theme-Wechsel ändert damit kein einziges Zeichen Inhalt.
%
%  Alles unten ist Blindtext einer fiktiven Person -- durch die eigenen
%  Angaben ersetzen. Die englische Fassung (cv-content-en.tex) hat dieselbe
%  Struktur; wer nur auf Deutsch schreibt, kann sie ignorieren.
%
%  MASTER-LEBENSLAUF. Diese Datei darf ALLE Angaben enthalten, die je in
%  einem Lebenslauf stehen sollen. Was eine Bewerbung davon zeigt,
%  entscheidet ihre meta.tex:
%      \cvshow{schluessel}    Baustein dazu, der sonst wegbleibt
%      \cvhide{schluessel}    Baustein weg, der sonst erscheint
%  "make cv-master" baut die vollständige Fassung zum Durchsehen.
%
%  In dieser Datei:
%      \cvopt{schluessel}{...}             nur mit \cvshow oder im Master
%      \cvstd{schluessel}{...}             immer, ausser mit \cvhide
%      \cvifopt{schluessel}{dann}{sonst}   zwei Fassungen desselben Textes
%
%  Schlüssel dieser Vorlage (Voreinstellung: aus, ausser markiert mit *):
%    photo *        Foto im Kopf
%    keywords *     Stichwortzeile unter dem Titel (kompakt, rueckgrat)
%    github         GitHub als weiterer Kontaktplatz
%    web            eigene Website als Link (rueckgrat)
%    nationality    Nationalität und Bewilligung als Kontaktplätze
%    exitreason *   Austrittsgrund in der Notizzeile der aktuellen Stelle
%    details        zwei zusätzliche Stichpunkte in der aktuellen Stelle
%    tools          Stichpunkt Berechnungswerkzeuge in der früheren Stelle
%    projects       Abschnitt Projekte nach der Berufserfahrung
%    grades         Abschlussnoten in der Ausbildung
%    further *      Abschnitt Weiterbildung
%    french *       Sprachen: Französisch
%    interests      Abschnitt Interessen
%    reference      benannte Referenz statt des Standardsatzes
%  Eigene Schlüssel entstehen einfach durch Verwenden: \cvopt{neu}{...}.
%
%  Weitere Schalter in der meta.tex:
%    \newcommand{\cvsummaryoverride}{...}   eigenes Kurzprofil für diese Bewerbung
%    \newcommand{\cvskillsoverride}{...}    eigene Kompetenzzeilen (\cvfact-Zeilen)
%    \newcommand{\cvheaderclassic}{}        kompakt: früherer Kopf, Kontakt dreispaltig
%    \newcommand{\noletters}{}  \newcommand{\nomotivation}{}  \newcommand{\lettersfirst}{}
%    \cvshow{coverpage}                     Deckblatt
%    \newcommand{\enclosures}{...}          Beilagenzeile im Anschreiben
%=========================================================================

%--- Stammdaten ----------------------------------------------------------

\newcommand{\CVname}{Max Muster}
% Titel unter dem Namen -- kommt pro Bewerbung aus meta.tex (\cvjobtitle)
\newcommand{\CVrole}{\cvjobtitle}
\newcommand{\CVphoto}{images/platzhalter.png}
% Bildausschnitt: \CVphotozoom = Bildbreite / Rahmenbreite (1 = bündig),
% \CVphotoshift verschiebt das Motiv vertikal (negativ = nach unten).
\renewcommand{\CVphotozoom}{1}
\renewcommand{\CVphotoshift}{0}
% Foto spiegeln, damit die Person zum Text hin blickt: -1 (Voreinstellung 1)
\renewcommand{\CVphotomirror}{1}
% Nur rueckgrat: Höhe des Fotos im Verhältnis zur Breite, 1.25 = Hochformat 4:5
\renewcommand{\CVphotoaspect}{1}
% Unterschrift als Bild über dem Namen im Brief, z.B. images/unterschrift.png (leer = keine)
\renewcommand{\CVsignature}{}
\newcommand{\CVaddress}{Musterstrasse 1, 8000 Zürich}
% Ort für die Datumszeile der Briefe
\newcommand{\CVcity}{Zürich}
\newcommand{\CVphone}{+41 79 000 00 00}
\newcommand{\CVmail}{max.muster@example.com}
\newcommand{\CVbirth}{01.01.1990}
\newcommand{\CVcivil}{ledig}
% Sechster Kontaktplatz: Symbol, Text und optional eine URL
\renewcommand{\CVlinkicon}{\icolinked}
\renewcommand{\CVlinktext}{max-muster}
\renewcommand{\CVlinkurl}{https://www.linkedin.com/in/max-muster/}
% Weitere Kontaktplätze, je über einen Schlüssel
\cvopt{github}{%
    \renewcommand{\CVgithubtext}{github.com/benutzername}%
    \renewcommand{\CVgithuburl}{https://github.com/}}
\cvopt{web}{%
    \renewcommand{\CVwebtext}{muster.example}%
    \renewcommand{\CVweburl}{https://www.example.com}}
\cvopt{nationality}{%
    \renewcommand{\CVnationality}{Schweizer}%
    \renewcommand{\CVpermit}{}}
\cvstd{keywords}{\renewcommand{\CVkeywords}{entwicklung · versuch · simulation · python}}

% Kopf im Layout kompakt: Kontaktangaben zweispaltig neben dem Foto
\ifdefined\cvheaderclassic\else\providecommand{\cvheadertwocol}{}\fi

%--- Kurzprofil ----------------------------------------------------------

\newcommand{\ContentSummary}{%
Ingenieur mit acht Jahren Erfahrung in Entwicklung und Erprobung mechatronischer
Baugruppen -- von der Auslegung über den Versuch bis zur Auswertung. Lorem ipsum dolor
sit amet, consetetur sadipscing elitr, sed diam nonumy eirmod tempor invidunt ut labore.%
}
% Eigenes Kurzprofil je Bewerbung: \newcommand{\cvsummaryoverride}{...} in der meta.tex
\ifdefined\cvsummaryoverride\renewcommand{\ContentSummary}{\cvsummaryoverride}\fi

%--- Berufserfahrung -----------------------------------------------------

\newcommand{\ContentExperience}{%
\begin{cvexperience}

\cvrole{Entwicklungsingenieur}{Muster Technik AG}{Zürich, Schweiz}{03/2020 -- heute}
\cvnote{03/2020 -- 12/2022 als Versuchsingenieur\cvstd{exitreason}{ • Austritt wegen Standortverlagerung}}
\begin{cvbullets}
    \item \lead{Auslegung:} Lorem ipsum dolor sit amet, consetetur sadipscing elitr,
    sed diam nonumy eirmod tempor invidunt ut labore et dolore magna aliquyam erat
    \item \lead{Versuch:} At vero eos et accusam et justo duo dolores et ea rebum --
    Stet clita kasd gubergren, no sea takimata sanctus est
    \item \lead{Auswertung:} Automatisierte Auswertung mit Python, lorem ipsum dolor
    sit amet, consetetur sadipscing elitr
    \cvopt{details}{\item \lead{Simulation:} Duis autem vel eum iriure dolor in hendrerit in
    vulputate velit esse molestie consequat
    \item \lead{Schulungen:} Vel illum dolore eu feugiat nulla facilisis at vero eros et
    accumsan}
\end{cvbullets}

\cvrole{Versuchsingenieur}{Beispiel GmbH}{Stuttgart, Deutschland}{10/2016 -- 02/2020}
\begin{cvbullets}
    \item \lead{Prüfstände:} Lorem ipsum dolor sit amet, consetetur sadipscing elitr,
    sed diam nonumy eirmod tempor invidunt ut labore
    \item \lead{Messtechnik:} Kraft-, Weg- und Temperaturmesstechnik, at vero eos et
    accusam et justo duo dolores
    \cvopt{tools}{\item \lead{Berechnungswerkzeuge:} Stet clita kasd gubergren, no sea
    takimata sanctus est lorem ipsum dolor sit amet}
\end{cvbullets}

% Zwei Zeiträume derselben Rolle beim selben Arbeitgeber: ein Eintrag.
\cvrole{Werkstudent}{Institut für Beispielforschung}{München, Deutschland}%
       {2015 -- 2016\cvalso{2013 -- 2014}}
\cvnote{Unterstützung in Forschungsprojekten zu Leichtbaustrukturen}

\end{cvexperience}%
}

%--- Projekte ------------------------------------------------------------
%  Abschnitt nach der Berufserfahrung, nur mit \cvshow{projects} oder im Master.

\renewcommand{\cvprojectsection}{\cvopt{projects}{%
    \cvsection[8]{\LblProjects}
    \begin{cvexperience}
    \cvrole{Werkstatt-Planer}{Web-Anwendung für eine offene Werkstatt}{\href{https://www.example.com}{example.com/werkstatt}}{2023 -- heute}
    \begin{cvbullets}
        \item \lead{Idee und Umsetzung:} Mitglieder reservieren Maschinen und sehen, wer eine
        Einführung geben kann
        \item \lead{Technik:} Python, Flask, SQLite
    \end{cvbullets}
    \end{cvexperience}%
}}

%--- Ausbildung ----------------------------------------------------------

\newcommand{\ContentEducation}{%
\begin{cvexperience}

\cvrole{Master of Science, Maschinenbau}{Technische Universität Musterstadt}{Musterstadt, Deutschland}{10/2014 -- 09/2016}
\cvnote{Schwerpunkt: Konstruktion und Simulation\cvopt{grades}{, Abschlussnote 1,5}}
\begin{cvbullets}
    \item \lead{Masterarbeit:} Lorem ipsum dolor sit amet, consetetur sadipscing elitr\cvopt{grades}{ (Note 1,3)}
\end{cvbullets}

\cvrole{Bachelor of Science, Maschinenbau}{Technische Universität Musterstadt}{Musterstadt, Deutschland}{10/2011 -- 09/2014}
\cvopt{grades}{\cvnote{Abschlussnote 2,0}}

\end{cvexperience}%
}

%--- Weiterbildung -------------------------------------------------------
%  Mit \cvhide{further} bleibt der Abschnitt samt Überschrift weg.

\newcommand{\ContentFurther}{}
\cvstd{further}{\renewcommand{\ContentFurther}{%
\begin{cvexperience}
\cvplain{Python für Ingenieure}{2021}{Lorem ipsum dolor sit amet, consetetur
sadipscing elitr}
\cvplain{Projektmanagement}{2019}{At vero eos et accusam et justo duo dolores}
\end{cvexperience}%
}}

%--- Kompetenzen ---------------------------------------------------------
%  Bewusst ohne Prozentbalken: eine Skala ohne Bezugsgrösse ist für den
%  Leser nicht überprüfbar und wird von Parsern als Leerraum gelesen.
%  Stattdessen Nennung dessen, was tatsächlich im Einsatz ist.

\newcommand{\ContentSkills}{%
\cvfact{Engineering}{FE-Simulation (ANSYS Workbench), CAD, Messtechnik}
\cvfact{Software}{Python, SQL, Git}
\cvfact{Systeme}{Siemens Teamcenter (PLM), SAP}
\cvfact{Methodik}{Versuchsplanung, FMEA, Reviews}%
}
% Eigene Kompetenzzeilen je Bewerbung: \newcommand{\cvskillsoverride}{\cvfact{..}{..}...}
% in der meta.tex. Der Master zeigt immer die vollständigen Zeilen.
\ifdefined\cvskillsoverride\ifdefined\cvmaster\else\renewcommand{\ContentSkills}{\cvskillsoverride}\fi\fi

% Kurzform als Schlagwort-Chips (nur das Sidebar-Theme zeigt sie)
\newcommand{\ContentTags}{%
\cvtag{Python}\cvtag{SQL}\cvtag{ANSYS}\cvtag{CAD}\cvtag{Messtechnik}%
\cvtag{Teamcenter}\cvtag{FMEA}%
}

%--- Sprachen ------------------------------------------------------------

\newcommand{\ContentLanguages}{%
\cvfact{Deutsch}{Muttersprache (C2)}
\cvfact{Englisch}{Verhandlungssicher (C1)}
\cvstd{french}{\cvfact{Französisch}{Grundkenntnisse (A2)}}%
}

%--- Zusatzabschnitte ----------------------------------------------------
%  \cvextrasections führt jedes Theme am Ende des Lebenslaufs aus -- für
%  Abschnitte wie Interessen oder Referenzen. Leer lassen = kein Abschnitt.

\renewcommand{\cvextrasections}{%
    \cvopt{interests}{%
        \cvsection{Interessen}
        \cvfact{Sport}{Velofahren, Skitouren}
        \cvfact{Werkstatt}{Holzarbeiten und 3D-Druck}}%
    \cvsection{\LblReferences}
    \cvifopt{reference}%
        {\cvfact{Erika Beispiel}{Leiterin Entwicklung, Muster Technik AG, Vorgesetzte}
         \cvfact{Kontakt}{+41 79 000 00 01}}%
        {\cvfact{}{Referenzen werden auf Nachfrage gerne mitgeteilt.}}%
}
